\section{Related Literature and Contribution Boundary}
\label{sec:literature}

The logic of coordinated industrialization predates the present framework. Rosenstein-Rodan's big-push argument formalized the possibility that individually unattractive investments can become viable when undertaken together. Hirschman's strategy of unbalanced growth placed particular emphasis on induced investment through backward and forward linkages. The present paper uses complementarity in a much narrower way: only as an entry gate for productive relations whose supply, demand, or input conditions are jointly activated. The fixed-point mathematics is inherited and is not itself a contribution.

The capability state in equation~\eqref{eq:capability-law} belongs to the broad learning-by-doing and technological-capability tradition. Arrow formalized learning-by-doing as an endogenous productivity mechanism. Pack and Westphal placed technological change and capability accumulation at the center of industrialization, and the literature on East Asian development has repeatedly emphasized the interaction among selective support, exports, imported technology, learning, and performance discipline. Westphal's account of Korea, for example, argues that selective intervention operated within an export-propelled system and contributed to the achievement of international competitiveness in a number of industries. Amsden and Wade likewise emphasize reciprocity, performance requirements, and the state's capacity to discipline supported firms rather than merely shelter them. These literatures already own the broad proposition that temporary support can be developmental when it generates learning rather than permanent dependence.

The playing-field extension in Section~\ref{sec:politics} is therefore not a claim to discover performance discipline. Its narrower purpose is to express a particular distinction already implicit in that literature: credibility that support will be honored during a learning window is not the same thing as credibility that the terminal condition cannot be privately renegotiated. The reduced-form functions $G(\sigma)$ and $D(\chi)$ isolate these two institutional margins and separate technological graduability from the incentive compatibility of pursuing graduation.

Modern industrial-policy work also emphasizes design and political economy rather than a simple intervention-versus-laissez-faire dichotomy. Harrison and Rodr\'iguez-Clare review both the theoretical rationales and the weak empirical basis for many hard interventions, while allowing an important role for softer cluster-level coordination. Aghion and coauthors find that sectoral policies can be more productivity-enhancing when they are competition-friendly. Liu shows how distortions and policy effects propagate through production networks. These contributions are important boundaries on the present exercise: a reduced-form graduation condition is not a welfare theorem and does not establish that a particular subsidy, tariff, or credit allocation is optimal.

The post-exit part of the paper sits near evolutionary economics, structural transformation, and the literature on related diversification. Schumpeter's creative destruction makes industrial replacement central to capitalist development. Nelson and Winter emphasize routines and the evolutionary persistence of firm capabilities. Hidalgo and coauthors show empirically that countries tend to move into products related to capabilities already present. Global-value-chain research studies upgrading, governance, and the relocation of production functions across firms and countries. The present term \emph{successful supersession} is therefore not a claim to have discovered that capabilities can migrate. It is a bookkeeping device used to distinguish capability-preserving exit from capability-destroying exit inside the same staged framework that begins with temporary support.

The positive-systems mathematics used below is also inherited. The coupled capability recursion is a standard nonnegative affine system, and the equivalence between stability and a nonsingular $M$-matrix condition is classical. The forward-invariance argument used for dynamic unsupported viability is likewise standard monotone-systems reasoning. The paper uses that machinery to expose implications of the economic state definitions; it does not claim the underlying mathematical tools as novel.

The narrowest potential contribution of the paper is consequently fourfold. First, it places an explicit capability state between supported activation and unsupported viability and derives a minimum support duration from that state law. Second, it evaluates graduation against recursively support-free dependencies and then separates instantaneous withdrawal survival from forward-invariant unsupported viability. Third, it formalizes a reduced-form playing-field condition distinguishing technological graduation from the relative private value of political dependence. Fourth, it makes the fate of capability after industry exit an explicit terminal diagnostic, allowing the same framework to distinguish failed graduation, destructive de-accretion, and successful supersession. The coupled-network result clarifies the mathematics connecting these pieces and identifies a partial bridge to reproductive accretion.
